\chapter{KẾT LUẬN}
\label{chap:conclusion}

Báo cáo đã khảo sát tám nghiên cứu mới nhất (tháng 5 -- tháng 9/2026) về multi-agent system và collaborative reasoning, phân loại chúng thành bốn hướng: safety \& robustness (MIRAGE, TIDE), orchestration \& control (ProgRouter, HalluProp), communication \& interaction (DAIN, AGENT-RADAR), và domain applications (MAC-RRG, UrbanAgent). Với mỗi nghiên cứu, báo cáo đã trình bày bài toán, phương pháp, thực nghiệm, kết quả và đánh giá phản biện; sau đó so sánh có hệ thống theo phương pháp, dữ liệu, kết quả, chi phí, ưu nhược điểm và phạm vi áp dụng.

Các nhận định chính gồm:
\begin{itemize}
  \item \textbf{Communication graph} là đối tượng trung tâm của lĩnh vực: nó vừa quyết định hiệu năng, vừa là kênh lan truyền rủi ro, vừa là tài sản cần bảo vệ.
  \item \textbf{Lan truyền} lỗi và thông tin sai là rủi ro cốt lõi; thảo luận nhiều vòng có xu hướng khuếch đại hơn là sửa sai, đi ngược giả định phổ biến rằng cộng tác luôn làm hệ thống chính xác hơn.
  \item \textbf{Hiệu quả} không đến từ việc dùng nhiều agent hơn hay mô hình lớn hơn, mà từ việc kích hoạt đúng agent, đúng mô hình và chú ý đúng nội dung.
  \item \textbf{Độ nghiêm ngặt thực nghiệm} của lĩnh vực còn hạn chế: quy mô nhỏ, thiếu kiểm định thống kê, thiếu đo chi phí và thiếu định nghĩa thống nhất về ``agent''.
\end{itemize}

Từ đó, báo cáo đề xuất các hướng phát triển: xây dựng benchmark đo đồng thời chất lượng -- chi phí -- robustness, thiết kế cơ chế giao tiếp có nhận thức nguồn gốc bằng chứng, kết hợp điều phối với dự báo rủi ro, phát triển phòng thủ chủ động trước thông tin sai và chuẩn hóa cách báo cáo kết quả \MAS. Nhóm cho rằng giai đoạn tiếp theo của lĩnh vực sẽ chuyển từ câu hỏi ``nhiều agent có tốt hơn một agent không'' sang câu hỏi ``làm thế nào để một multi-agent system vừa hiệu quả, vừa đáng tin cậy và vừa an toàn''.
